Let
\[
 \widetilde\tau=\frac1n\sum_{i=1}^n Z_i,
 \qquad
 Z_i=\sum_{a\in S_c}
 \frac{c_a\mathbf 1\{A_i=a\}}{q_a^\star}
 \bigl(Y_i^{\mathrm{obs}}-1/2\bigr).
\]
Because \(c\neq0\), one has \(L_c>0\).  For every \(a\in S_c\), the definition \(q_a^\star=|c_a|/L_c\) therefore gives \(q_a^\star>0\), so all displayed divisions are valid.  Under the independent design in \(\mathrm{def:first\text{-}order\text{-}procedure}\), and for a fixed response schedule \(z\),
\[
 \begin{aligned}
 \mathbb E_{\mathcal D}(Z_i\mid z)
 &=\sum_{a\in S_c}c_a(t_{i,a}-1/2)\\
 &=\sum_{a\in\mathcal A_K}c_at_{i,a}=:\mu_i.
 \end{aligned}
\]
The second equality uses \(c_a=0\) off \(S_c\) and \(\sum_a c_a=0\).  Hence
\[
 \mathbb E_{\mathcal D}(\widetilde\tau\mid z)
 =\frac1n\sum_{i=1}^n\mu_i=\tau_c(z).
\]
Since \(t_{i,a}\in\{0,1\}\), the identity \((t_{i,a}-1/2)^2=1/4\) gives
\[
 \begin{aligned}
 \mathbb E_{\mathcal D}(Z_i^2\mid z)
 &=\frac14\sum_{a\in S_c}\frac{c_a^2}{q_a^\star}
 =\frac14\sum_{a\in S_c}L_c|c_a|\\
 &=\frac{L_c^2}{4}=C_0(c).
 \end{aligned}
\]
Assignments are independent across units, so the variables \(Z_i\) are independent conditional on the fixed schedule.  Therefore
\[
 \begin{aligned}
 R_n(\mathcal D,\widetilde\tau;z)
 &=\frac1{n^2}\sum_{i=1}^n
 \operatorname{Var}_{\mathcal D}(Z_i\mid z)\\
 &=\frac1{n^2}\sum_{i=1}^n\{C_0(c)-\mu_i^2\}
 \leq\frac{C_0(c)}n.
 \end{aligned}
\]

Write \(P_c=\{a:c_a>0\}\) and \(N_c=\{a:c_a<0\}\).  The zero-sum identity implies
\[
 \sum_{a\in P_c}c_a
 =-\sum_{a\in N_c}c_a=\frac{L_c}{2}.
\]
For every binary response vector \(t_i\), it follows that
\[
 -\frac{L_c}{2}\leq\sum_a c_at_{i,a}\leq\frac{L_c}{2},
\]
and averaging over units yields \(\tau_c(z)\in[-L_c/2,L_c/2]\).  Metric projection onto a closed interval cannot increase squared distance from any point of that interval.  Thus the projected estimator \(\widehat\tau_{\mathrm{cHT}}^\star\) belongs to the stipulated estimator action space and satisfies
\[
 \max_{z\in\mathcal T_K^n}
 R_n(\mathcal D,\widehat\tau_{\mathrm{cHT}}^\star;z)
 \leq\frac{C_0(c)}n.
\]
In particular, \(\rho_n(c)\leq C_0(c)/n\).

We next retain a direct first-order converse based only on a prior over fixed binary schedules.  Since \(K\) is fixed, \(s:=|S_c|\) is fixed.  Put \(\epsilon_n=n^{-1/4}/4\) and
\[
 w(u)=\frac{15}{16}(1-u^2)^2\mathbf 1\{|u|\leq1\},
 \qquad
 J=\int_{-1}^1\frac{\{w'(u)\}^2}{w(u)}\,du<\infty.
\]
For each \(a\in S_c\), independently draw \(P_a\) with density
\[
 p\longmapsto\epsilon_n^{-1}w\!\left(\frac{p-1/2}{\epsilon_n}\right),
\]
and set \(P_a=1/2\) for \(a\notin S_c\).  Independently draw \(U_1,\ldots,U_n\) identically and uniformly on \([0,1]\), and define
\[
 Y_i(a)=\mathbf 1\{U_i\leq P_a\}.
\]
After \((P,U)\) is drawn, this construction is a fixed element of \(\mathcal T_K^n\); thus it is a legitimate prior on the finite schedule space and is independent of every admissible assignment law.

Define
\[
 \vartheta(P)=\sum_{a\in S_c}c_aP_a,
 \qquad
 g_P(u)=\sum_{a\in S_c}c_a\mathbf 1\{u\leq P_a\}.
\]
Conditional on \(P\),
\[
 \tau_c(z)=\frac1n\sum_{i=1}^ng_P(U_i),
 \qquad
 \mathbb E\{\tau_c(z)\mid P\}=\vartheta(P).
\]
The zero-sum identity makes \(g_P\) vanish outside
\([\min_{a\in S_c}P_a,\max_{a\in S_c}P_a]\), an interval of length at most \(2\epsilon_n\).  Also \(|g_P(u)|\leq L_c/2\).  Hence
\[
 \begin{aligned}
 e_n^2
 &:=\mathbb E\{\tau_c(z)-\vartheta(P)\}^2\\
 &=\frac1n\mathbb E\{\operatorname{Var}(g_P(U_1)\mid P)\}
 \leq\frac{L_c^2\epsilon_n}{2n},
 \end{aligned}
\]
so \(e_n=O(n^{-5/8})=o(n^{-1/2})\).

Fix an arbitrary assignment law \(\mathcal D\), allowing arbitrary dependence among the assignments, and put \(s_a=\mathbb E_{\mathcal D}r_a\).  Conditional on \((P,A)\), the observed outcomes are independent and satisfy
\[
 Y_i^{\mathrm{obs}}\sim\operatorname{Bernoulli}(P_{A_i}).
\]
By the symmetry and compact support of \(w\),
\[
 \kappa_n
 :=\mathbb E\!\left\{\frac1{P_a(1-P_a)}\right\}
 =4+O(\epsilon_n^2).
\]
For \(a\in S_c\), let \(S_a\) be the score of the joint prior--likelihood density with respect to \(P_a\).  Conditional mean-zero of each Bernoulli likelihood score, independence of the arm priors, and independence of the observed outcomes conditional on \((P,A)\) show that the score covariance matrix is diagonal, with
\[
 I_a:=\mathbb E(S_a^2)
 =\kappa_ns_a+\frac{J}{\epsilon_n^2}>0.
\]
Thus no exposure-positivity condition is required for this converse; the prior-information term makes every denominator positive even when \(s_a=0\).

Let \(T=T(A,Y^{\mathrm{obs}})\) be any possibly biased estimator.  The functions \(w\) and \(w'\) vanish at both endpoints of their support.  Integration by parts in \(P_a\), with no boundary term, gives
\[
 \mathbb E(TS_a)=0,
 \qquad
 \mathbb E\{\vartheta(P)S_a\}=-c_a,
 \qquad
 \mathbb E[\{T-\vartheta(P)\}S_a]=c_a.
\]
For every vector \(v=(v_a)_{a\in S_c}\), Cauchy--Schwarz and diagonality of the score covariance give
\[
 \mathbb E\{T-\vartheta(P)\}^2
 \geq\frac{(v^\top c)^2}{\sum_{a\in S_c}v_a^2I_a}.
\]
Taking \(v_a=c_a/I_a\), and then applying Cauchy--Schwarz to the vectors \((|c_a|/\sqrt{I_a})_a\) and \((\sqrt{I_a})_a\), yields
\[
 \begin{aligned}
 \mathbb E\{T-\vartheta(P)\}^2
 &\geq\sum_{a\in S_c}\frac{c_a^2}{I_a}
 \geq\frac{L_c^2}{\sum_{a\in S_c}I_a}\\
 &\geq\frac{L_c^2}{\kappa_n n+sJ/\epsilon_n^2}
 =\frac{C_0(c)}n+o(n^{-1}).
 \end{aligned}
\]
Here \(\sum_{a\in S_c}s_a\leq n\), \(\kappa_n=4+O(n^{-1/2})\), and \(sJ/\epsilon_n^2=O(n^{1/2})\).

The reverse triangle inequality in \(L^2\) now gives, uniformly over \(\mathcal D\) and \(T\),
\[
 \left[\mathbb E\{T-\tau_c(z)\}^2\right]^{1/2}
 \geq
 \frac{L_c}{\sqrt{\kappa_n n+sJ/\epsilon_n^2}}-e_n.
\]
The first term is asymptotic to \(L_c/(2\sqrt n)\), while \(e_n=o(n^{-1/2})\).  Squaring the positive part therefore proves
\[
 \mathbb E\{T-\tau_c(z)\}^2
 \geq\frac{C_0(c)}n+o(n^{-1}).
\]
The maximum over fixed schedules dominates the average under the constructed prior.  Taking the infimum over all assignment laws and all possibly biased estimators proves the stated unrestricted first-order lower bound.

There is also a sharper finite-sample converse.  By \(\mathrm{thm:embedded\text{-}two\text{-}arm\text{-}converse}\), for every \(n\geq1\),
\[
 \rho_n(c)\geq C_0(c)\rho_{n,2}.
\]
Its published two-arm consequence is
\[
 \rho_n(c)
 \geq\frac{C_0(c)}n-C_0(c)C_A n^{-4/3}+o(n^{-4/3}),
\]
which in particular recovers the preceding \(C_0(c)/n+o(n^{-1})\) converse.  When \(|S_c|=2\), the embedded theorem states equality.  When \(|S_c|\geq3\), only the displayed inequality is used: no second-order equality is claimed.

Combining either first-order converse with the finite-sample upper certificate yields
\[
 \lim_{n\to\infty}n\rho_n(c)
 =C_0(c)=\frac{L_c^2}{4}
 =\frac14\left(\sum_{a\in\mathcal A_K}|c_a|\right)^2.
\]
The equality \(\rho_n(c)=G_n(c)\) from \(\mathrm{thm:exact\text{-}response\text{-}type\text{-}game}\) transfers the same first-order value and the finite-sample certificates to the explicit orbit game.

Finally, the allocation choice has a matching convex optimality certificate.  For any allocation vector \(q\) with \(q_a>0\) on \(S_c\), Cauchy--Schwarz gives
\[
 \frac14\sum_{a\in S_c}\frac{c_a^2}{q_a}
 \geq\frac14
 \frac{\left(\sum_{a\in S_c}|c_a|\right)^2}
 {\sum_{a\in S_c}q_a}
 \geq\frac{L_c^2}{4}.
\]
Equality holds precisely when no probability is assigned outside \(S_c\) and \(q_a\) is proportional to \(|c_a|\) on \(S_c\), namely \(q=q^\star\).  This completes the first-order saddle and allocation-optimality proof.